\subsection{معرفی و انگیزه}
این مکانیزم بر مبنای مبانی نظری پیشرفته در ادبیات توجه و فشرده‌سازی \cite{ref_lolaattention_2024} توسعه یافته است.

مکانیزم توجه \lr{LoLA (Low-Rank Linear Attention with Sparse Caching)} به منظور برطرف ساختن دو چالش عمده در استقرار مدل‌های ترنسفورمر مقیاس‌بزرگ طراحی شده است: نخست، حجم عظیم حافظه مصرفی ناشی از ماتریس‌های پروجکشن کامل، و دوم، رشد نامحدود و خطی بافر تانسورهای پیشین (\lr{KV-Cache}) در دیکودرهای خودبازگشتی. در کاربردهای بلادرنگ و سامانه‌های لبه (\lr{Edge Devices})، سربار ناشی از جابه‌جایی داده‌های حافظه پنهان بین حافظه اصلی و پردازنده گرافیکی عامل اصلی افت توان پردازشی به شمار می‌رود.

نوآوری ماژول \lr{LoLA} در ترکیب دو راهبرد مکمل است: از یک سو، استفاده از نگاشت‌های خطی کم‌رتبه (\lr{LowRankLinear}) مبتنی بر تجزیه مقادیر منفرد (\lr{SVD}) به همراه اتصالات باقیمانده یادگیرنده رتبه-پایین (\lr{LoRA}) برای کاهش وزن‌های پایه؛ و از سوی دیگر، طراحی یک سامانه سلسله‌مراتبی کش چندبخشی (\lr{LowRankCache}) که توکن‌های اخیر را به صورت دقیق، توکن‌های تاریخی را با فواصل گزینشی (\lr{Strided/Sparse})، و توکن‌های دوردست را به صورت زیرفضای تصویرشده ارتوگونال با رتبه تقلیل‌یافته (\lr{Orthogonal Low-Rank QR Projection}) نگهداری می‌نماید.
